% =====================================================================
\chapter{Regresión lineal simple y correlación}
\label{tema:6}
% =====================================================================

La regresión estudia cómo cambia en promedio una variable respuesta $y$ cuando cambia una variable explicativa $x$. Este tema deduce la recta de mínimos cuadrados, estudia los residuos y la descomposición de la variabilidad que da lugar al coeficiente de determinación $R^{2}$, muestra cómo las transformaciones logarítmicas permiten ajustar leyes de potencia y termina con la distinción entre correlación y causalidad. Se sigue el capítulo~12 de OpenStax~\cite{openstax} (§12.3 «The Regression Equation»), el capítulo~14 de Rice~\cite{rice}, «Linear Least Squares», el capítulo~13 de Wasserman~\cite{wasserman}, «Linear and Logistic Regression», y la lectura~26 del curso 18.05 del MIT~\cite{mit1805}; para la causalidad, el capítulo~16 de Wasserman, «Causal Inference».

\section{La recta de mínimos cuadrados}

Se dispone de $n\ge2$ pares $(x_1,y_1),\dots,(x_n,y_n)$ con los $x_i$ no todos iguales. Se usan las sumas centradas
\[
S_{xx}=\sum_{i=1}^{n}(x_i-\xbar)^{2},\qquad S_{yy}=\sum_{i=1}^{n}(y_i-\ybar)^{2},\qquad S_{xy}=\sum_{i=1}^{n}(x_i-\xbar)(y_i-\ybar),
\]
de modo que $s_x^{2}=S_{xx}/(n-1)$, $s_{xy}=S_{xy}/(n-1)$ y $r=S_{xy}/\sqrt{S_{xx}S_{yy}}$ (\cref{def:t1-cov}). Por la \cref{prop:t1-abreviada}, $S_{xx}=\sum x_i^{2}-n\xbar^{2}$ y, análogamente, $S_{xy}=\sum x_iy_i-n\xbar\ybar$.

\begin{definicion}[Recta de mínimos cuadrados]\label{def:t6-mc}
Para cada recta $y=b_0+b_1x$, el \emph{valor ajustado} del dato $i$ es $\widehat y_i=b_0+b_1x_i$ y el \emph{residuo} es $e_i=y_i-\widehat y_i$. La \emph{recta de mínimos cuadrados} (o de regresión de $y$ sobre $x$) es la que minimiza la suma de cuadrados de los residuos
\[
\operatorname{SC}(b_0,b_1)=\sum_{i=1}^{n}(y_i-b_0-b_1x_i)^{2}.
\]
\end{definicion}

\begin{teorema}[Coeficientes de mínimos cuadrados]\label{teo:t6-mc}
Si $S_{xx}>0$, la función $\operatorname{SC}$ tiene un único mínimo, en
\[
\widehat b_1=\frac{S_{xy}}{S_{xx}}=r\,\frac{s_y}{s_x},\qquad\widehat b_0=\ybar-\widehat b_1\xbar,
\]
y el valor mínimo es $S_{yy}-S_{xy}^{2}/S_{xx}=S_{yy}(1-r^{2})$.
\end{teorema}

\begin{proof}
\emph{Paso 1 (separar la ordenada).} Sea $c=\ybar-b_0-b_1\xbar$. Entonces $y_i-b_0-b_1x_i=\bigl[(y_i-\ybar)-b_1(x_i-\xbar)\bigr]+c$. Al desarrollar el cuadrado y sumar, el término cruzado es $2c\sum_i\bigl[(y_i-\ybar)-b_1(x_i-\xbar)\bigr]=0$, porque las desviaciones respecto de la media suman cero (\cref{prop:t1-extremal}). Por tanto,
\[
\operatorname{SC}(b_0,b_1)=\sum_{i}\bigl[(y_i-\ybar)-b_1(x_i-\xbar)\bigr]^{2}+n\,c^{2}.
\]
\emph{Paso 2 (completar el cuadrado en $b_1$).} Desarrollando el primer sumando,
\[
\sum_i\bigl[(y_i-\ybar)-b_1(x_i-\xbar)\bigr]^{2}=S_{yy}-2b_1S_{xy}+b_1^{2}S_{xx}=S_{xx}\Bigl(b_1-\frac{S_{xy}}{S_{xx}}\Bigr)^{2}+S_{yy}-\frac{S_{xy}^{2}}{S_{xx}}.
\]
\emph{Paso 3 (minimizar).} Como $S_{xx}>0$, $\operatorname{SC}(b_0,b_1)\ge S_{yy}-S_{xy}^{2}/S_{xx}$, con igualdad si y solo si $b_1=S_{xy}/S_{xx}$ y $c=0$, es decir, $b_0=\ybar-b_1\xbar$. Finalmente, $S_{xy}/S_{xx}=r\sqrt{S_{yy}/S_{xx}}=r\,s_y/s_x$ y $S_{xy}^{2}/S_{xx}=r^{2}S_{yy}$.
\end{proof}

La condición $\widehat b_0=\ybar-\widehat b_1\xbar$ dice que la recta pasa por el punto de medias $(\xbar,\ybar)$. La pendiente tiene el signo de $r$: es positiva cuando la asociación es creciente.

\begin{ejemplo}\label{ej:t6-recta}
Para los datos del \cref{ej:t1-r}, $(1,2),(2,4),(3,5),(4,4),(5,5)$: $S_{xx}=10$, $S_{yy}=6$ y $S_{xy}=6$. Entonces $\widehat b_1=0{,}6$, $\widehat b_0=4-0{,}6\cdot3=2{,}2$ y la recta es $\widehat y=2{,}2+0{,}6x$ (\cref{fig:t6-recta}). Los residuos son $-0{,}8;\allowbreak\ 0{,}6;\allowbreak\ 1{,}0;\allowbreak\ -0{,}6;\allowbreak\ -0{,}2$, que suman cero, y su suma de cuadrados es $2{,}4=6\,(1-0{,}6)$.
\end{ejemplo}

\begin{figure}[ht]
\centering
\begin{tikzpicture}[x=1.1cm,y=0.75cm]
  \draw[->] (0,0) -- (6,0) node[right]{\small $x$};
  \draw[->] (0,0) -- (0,6.2) node[above]{\small $y$};
  \foreach \i in {1,...,5}{\draw (\i,0) -- (\i,-0.12) node[below]{\scriptsize\i};}
  \foreach \j in {1,...,6}{\draw (0,\j) -- (-0.1,\j) node[left]{\scriptsize\j};}
  \draw[thick] (0.5,2.5) -- (5.6,5.56) node[right]{\small $\widehat y=2{,}2+0{,}6x$};
  \foreach \x/\y/\yh in {1/2/2.8,2/4/3.4,3/5/4.0,4/4/4.6,5/5/5.2}{
    \draw[dashed,gray] (\x,\y) -- (\x,\yh);
    \fill (\x,\y) circle (2pt);
  }
  \draw (3,4) circle (3pt);
  \node[below right] at (3,4) {\scriptsize $(\xbar,\ybar)$};
\end{tikzpicture}
\caption{Recta de mínimos cuadrados del \cref{ej:t6-recta}. Los segmentos discontinuos son los residuos; la recta pasa por el punto de medias $(3,4)$.}\label{fig:t6-recta}
\end{figure}

\section{Residuos y coeficiente de determinación}

\begin{proposicion}[Propiedades de los residuos]\label{prop:t6-residuos}
Para la recta de mínimos cuadrados, $\sum_ie_i=0$, $\sum_ix_ie_i=0$ y $\sum_i\widehat y_ie_i=0$.
\end{proposicion}

La demostración es el \cref{eje:t6-residuos}. Las dos primeras igualdades son las \emph{ecuaciones normales} (\cref{eje:t6-normales}): los residuos son ortogonales, como vectores de $\R^{n}$, al vector de unos y al vector de las $x_i$. Este es el punto de vista geométrico que el módulo~8 generaliza a varias variables explicativas.

\begin{teorema}[Descomposición de la variabilidad]\label{teo:t6-anova}
Sean $\operatorname{SC}_T=\sum_i(y_i-\ybar)^{2}=S_{yy}$ (suma de cuadrados total), $\operatorname{SC}_R=\sum_i(\widehat y_i-\ybar)^{2}$ (explicada por la regresión) y $\operatorname{SC}_E=\sum_ie_i^{2}$ (residual o de error). Entonces
\[
\operatorname{SC}_T=\operatorname{SC}_R+\operatorname{SC}_E.
\]
\end{teorema}

\begin{proof}
Se escribe $y_i-\ybar=(\widehat y_i-\ybar)+e_i$ y se eleva al cuadrado:
$\sum_i(y_i-\ybar)^{2}=\sum_i(\widehat y_i-\ybar)^{2}+\sum_ie_i^{2}+2\sum_i(\widehat y_i-\ybar)e_i$.
El último término es $2\sum_i\widehat y_ie_i-2\ybar\sum_ie_i=0-0$ por la \cref{prop:t6-residuos}.
\end{proof}

\begin{definicion}[Coeficiente de determinación]\label{def:t6-r2}
El \emph{coeficiente de determinación} es $R^{2}=\operatorname{SC}_R/\operatorname{SC}_T=1-\operatorname{SC}_E/\operatorname{SC}_T$, la fracción de la variabilidad de $y$ explicada por la recta. Por el \cref{teo:t6-anova}, $0\le R^{2}\le1$.
\end{definicion}

\begin{corolario}\label{cor:t6-r2}
En la regresión lineal simple, $R^{2}=r^{2}$.
\end{corolario}

\begin{proof}
Por el \cref{teo:t6-mc}, $\operatorname{SC}_E=S_{yy}(1-r^{2})$, luego $R^{2}=1-\operatorname{SC}_E/S_{yy}=r^{2}$.
\end{proof}

En el \cref{ej:t6-recta}, $R^{2}=0{,}6$: la recta explica el $60\,\%$ de la variabilidad de $y$. La dispersión de los puntos alrededor de la recta se resume en el \emph{error típico residual} $s_e=\sqrt{\operatorname{SC}_E/(n-2)}$; se divide por $n-2$ porque se han estimado dos parámetros (compárese con la \cref{obs:t1-n-1}).

\section{Diagnóstico con residuos}

Un $R^{2}$ alto no garantiza que el modelo lineal sea adecuado. La herramienta básica de diagnóstico es el \emph{gráfico de residuos}, que representa $e_i$ frente a $x_i$ (o frente a $\widehat y_i$). Si el modelo es adecuado, los residuos se distribuyen sin estructura alrededor de cero. Algunas señales de alarma~\cite[§12.3]{openstax}:
\begin{enumerate}
  \item una curvatura sistemática (residuos positivos en los extremos y negativos en el centro, o al revés) indica que la relación no es lineal;
  \item una dispersión que crece con $x$ (forma de embudo) indica varianza no constante;
  \item un punto con un residuo muy grande es un posible atípico, y un punto con $x_i$ muy alejado de $\xbar$ puede ser muy \emph{influyente}: puede cambiar mucho la recta por sí solo.
\end{enumerate}
Además, la recta solo describe la relación en el rango de los $x_i$ observados: \emph{extrapolar} fuera de ese rango (\cref{eje:t6-buceo}) es arriesgado.

\begin{observacion}[El modelo estadístico]\label{obs:t6-modelo}
Si se supone que $y_i=\beta_0+\beta_1x_i+\varepsilon_i$ con errores $\varepsilon_i\sim\Nor(0,\sigma^{2})$ \iid, la log-verosimilitud de $(\beta_0,\beta_1)$ es una constante menos $\frac{1}{2\sigma^{2}}\operatorname{SC}(\beta_0,\beta_1)$; por tanto, los estimadores de mínimos cuadrados son los estimadores de máxima verosimilitud~\cite[cap.~13]{wasserman}, \cite[\emph{problem set}~6]{mit1805}. Con este modelo pueden construirse intervalos de confianza y contrastes para la pendiente (ampliación), que se estudian en~\cite[cap.~14]{rice}.
\end{observacion}

\section{Transformaciones: ajuste de leyes de potencia}

Muchas relaciones no lineales se vuelven lineales tras una transformación. Una \emph{ley de potencia} $y=a\,x^{b}$ con $a,x,y>0$ cumple
\[
\log y=\log a+b\log x,
\]
de modo que en escala log-log es una recta de pendiente $b$ y ordenada $\log a$. Para ajustarla se aplica mínimos cuadrados a los pares $(\log x_i,\log y_i)$ y se deshace la transformación: $\widehat a=10^{\widehat b_0}$ (o $e^{\widehat b_0}$ si se usan logaritmos naturales) y $\widehat b=\widehat b_1$. Análogamente, una relación exponencial $y=a\,e^{bx}$ se linealiza con $\ln y=\ln a+bx$. Obsérvese que así se minimizan los errores \emph{relativos} (en escala logarítmica), no los absolutos.

\begin{ejemplo}\label{ej:t6-potencia}
Si la pérdida de un modelo cumpliera exactamente $L=20\,N^{-0{,}1}$, al multiplicar el número de parámetros $N$ por $10$ la pérdida se multiplicaría por $10^{-0{,}1}\approx0{,}794$, y en escala log-log los puntos estarían sobre una recta de pendiente $-0{,}1$. El \cref{eje:t6-loglog} ajusta una ley así a partir de cuatro mediciones.
\end{ejemplo}

\section{Correlación y causalidad}

Una correlación entre $x$ e $y$ puede deberse a que $x$ influye en $y$, a que $y$ influye en $x$ (causalidad inversa), a que una tercera variable $z$ influye en ambas (\emph{variable de confusión}) o simplemente al azar. Los datos observacionales, por sí solos, no permiten distinguir estos casos~\cite[cap.~16]{wasserman}.

\begin{ejemplo}[Variable de confusión]\label{ej:t6-confusion}
En un conjunto de modelos de lenguaje, la longitud media de las respuestas y la puntuación en un benchmark de razonamiento pueden estar fuertemente correlacionadas. De ahí no se sigue que alargar artificialmente las respuestas mejore la puntuación: ambos rasgos pueden deberse a una tercera variable, como el tamaño del modelo o el entrenamiento con cadenas de razonamiento.
\end{ejemplo}

Un caso extremo de confusión es la \emph{paradoja de Simpson}: una asociación que se da en cada uno de los grupos de una población puede invertirse al agregar los grupos, cuando los tamaños de los grupos difieren entre las condiciones comparadas (\cref{eje:t6-simpson}).

La forma más fiable de establecer una relación causal es el \emph{experimento aleatorizado}: si el tratamiento (por ejemplo, el \emph{prompt} A o el B) se asigna al azar a las unidades (las preguntas), las variables de confusión se reparten en promedio por igual entre los grupos, y una diferencia sistemática en la respuesta solo puede atribuirse al tratamiento o al azar, cuya magnitud se cuantifica con los contrastes del \cref{tema:5}.

\begin{observacion}[Errores frecuentes]\label{obs:t6-errores}
(a) Interpretar la pendiente como efecto causal con datos observacionales.
(b) Extrapolar la recta fuera del rango de los datos.
(c) Concluir que no hay relación porque $r\approx0$ (la relación puede ser no lineal; \cref{eje:t6-cuadratica}).
(d) Creer que la recta de regresión de $x$ sobre $y$ es la inversa de la de $y$ sobre $x$: solo lo es si $\abs r=1$ (\cref{eje:t6-resumen}).
(e) Juzgar un ajuste solo por $R^{2}$, sin mirar el gráfico de residuos.
\end{observacion}

\begin{observacion}[Conexión con el curso]
El módulo~8 generaliza la recta de mínimos cuadrados a varias variables mediante las ecuaciones normales $X^{\top}X\widehat{\bm\beta}=X^{\top}\mathbf y$ y su interpretación como proyección, de la que la \cref{prop:t6-residuos} es el caso más sencillo, y demuestra que mínimos cuadrados es máxima verosimilitud con errores normales (\cref{obs:t6-modelo}). El módulo~16 ajusta en escala log-log las leyes de escala que relacionan la pérdida con el número de parámetros, de datos y de operaciones (\cref{ej:t6-potencia}). En el módulo~23, las sondas lineales detectan \emph{correlaciones} entre activaciones y conceptos, mientras que las intervenciones causales (modificar una activación y observar el efecto) son el análogo del experimento aleatorizado.
\end{observacion}

\begin{lecturas}
\item OpenStax~\cite{openstax}, §12.3 «The Regression Equation» y §12.4 «Testing the Significance of the Correlation Coefficient».
\item Rice~\cite{rice}, cap.~14 «Linear Least Squares».
\item Wasserman~\cite{wasserman}, cap.~13 «Linear and Logistic Regression» y cap.~16 «Causal Inference».
\item Orloff y Bloom~\cite{mit1805}, lectura~26 «Linear regression».
\item Stats~200~\cite{stats200}, clase~25 «The linear model».
\end{lecturas}

\section*{Ejercicios del tema 6}
\addcontentsline{toc}{section}{Ejercicios del tema 6}

\begin{ejercicio}{1}{eje:t6-vf}
Indique si cada afirmación es verdadera o falsa. (a) La recta de mínimos cuadrados pasa siempre por $(\xbar,\ybar)$. (b) Si $r=-0{,}9$, entonces $R^{2}=-0{,}81$. (c) Un $R^{2}$ cercano a~$1$ prueba que $x$ causa $y$. (d) La suma de los residuos de mínimos cuadrados es cero. (e) Si se intercambian los papeles de $x$ e $y$, la pendiente pasa a ser $1/\widehat b_1$.
\end{ejercicio}

\begin{ejercicio}{2}{eje:t6-basico}
Para los datos $(0,1),(1,3),(2,2),(3,5),(4,4)$: (a) calcule $S_{xx}$, $S_{yy}$, $S_{xy}$ y $r$; (b) obtenga la recta de mínimos cuadrados; (c) calcule los residuos, $\operatorname{SC}_E$, $R^{2}$ y $s_e$; (d) prediga $y$ para $x=2{,}5$.
\end{ejercicio}

\begin{ejercicio}{2}{eje:t6-gauss}
Se ajusta el modelo $y_i=ax_i+b+\varepsilon_i$ con errores $\Nor(0,\sigma^{2})$ \iid\ a los datos $(1,8),(3,2),(5,1)$. (a) Escriba la log-verosimilitud de un dato $(x_1,y_1)$. (b) Obtenga los EMV de $a$ y $b$ y compruebe que no dependen de $\sigma$.
\fuente{adaptado de Orloff y Bloom, MIT 18.05~\cite{mit1805}, \emph{problem set}~6, problema~2.}
\end{ejercicio}

\begin{ejercicio}{2}{eje:t6-examenes}
Las notas de once estudiantes en un tercer examen ($x$, sobre $80$) y en el final ($y$, sobre $200$) cumplen $n=11$, $\sum x_i=761$, $\sum y_i=1765$, $\sum x_i^{2}=52\,729$, $\sum y_i^{2}=287\,529$ y $\sum x_iy_i=122\,500$. (a) Calcule la recta de regresión y $r$. (b) Prediga la nota final de un estudiante con $73$ en el tercer examen. (c) Un estudiante tiene $(65,175)$; calcule su residuo.
\fuente{adaptado de OpenStax, \emph{Introductory Statistics}~2e~\cite{openstax}, §12.3, ejemplo~12.6.}
\end{ejercicio}

\begin{ejercicio}{2}{eje:t6-buceo}
La profundidad de una inmersión ($x$, en pies) y el tiempo máximo de inmersión ($y$, en minutos) son $(50,80),(60,55),(70,45),(80,35),(90,25),(100,22)$. (a) Calcule la recta de regresión y $r$. (b) Prediga $y$ para $x=85$ y para $x=110$. ¿Cuál de las dos predicciones es fiable? (c) Observe el signo de los residuos en los extremos y en el centro: ¿qué sugiere?
\fuente{adaptado de OpenStax, \emph{Introductory Statistics}~2e~\cite{openstax}, §12.3, \emph{Try It}~12.6.}
\end{ejercicio}

\begin{ejercicio}{2}{eje:t6-helados}
La temperatura máxima ($x$, en grados Fahrenheit) y el número de helados vendidos ($y$) en nueve días son $(70,105)$, $(85,240)$, $(65,49)$, $(72,147)$, $(80,231)$, $(61,38)$, $(75,193)$, $(78,196)$, $(68,89)$. Se dan $\sum x_i=654$, $\sum y_i=1288$, $\sum x_i^{2}=47\,988$, $\sum y_i^{2}=231\,026$ y $\sum x_iy_i=98\,132$. Calcule la recta de regresión, $r$ y $R^{2}$, y prediga las ventas de un día a $74$ grados.
\fuente{adaptado de OpenStax, \emph{Introductory Statistics}~2e~\cite{openstax}, §12.3, ejemplo «Ice Cream Cones».}
\end{ejercicio}

\begin{ejercicio}{2}{eje:t6-resumen}
De un conjunto de pares se conocen $\xbar=10$, $\ybar=50$, $s_x=2$, $s_y=8$ y $r=0{,}75$. (a) Obtenga la recta de regresión de $y$ sobre $x$ y prediga $y$ para $x=12$. (b) Obtenga la recta de regresión de $x$ sobre $y$. (c) Compruebe que el producto de las dos pendientes es $r^{2}$ y que las dos rectas no coinciden.
\end{ejercicio}

\begin{ejercicio}{2}{eje:t6-loglog}
Se mide la pérdida $L$ de cuatro modelos con $N=10^{6},10^{7},10^{8},10^{9}$ parámetros y se obtiene $L=5{,}0;\allowbreak\ 4{,}1;\allowbreak\ 3{,}3;\allowbreak\ 2{,}7$, cuyos logaritmos decimales son $0{,}6990;\allowbreak\ 0{,}6128;\allowbreak\ 0{,}5185;\allowbreak\ 0{,}4314$. (a) Ajuste una ley de potencia $L=aN^{-\alpha}$ por mínimos cuadrados en escala log-log. (b) Prediga la pérdida para $N=10^{10}$. (c) ¿Por qué esta predicción es una extrapolación?
\end{ejercicio}

\begin{ejercicio}{2}{eje:t6-simpson}
Dos modelos se evalúan en un banco con preguntas fáciles y difíciles. El modelo A acierta $45$ de $50$ fáciles y $75$ de $250$ difíciles; el modelo B acierta $200$ de $250$ fáciles y $10$ de $50$ difíciles. (a) Calcule las exactitudes de cada modelo en cada tipo de pregunta y en total. (b) ¿Qué modelo es mejor? Explique la aparente contradicción.
\end{ejercicio}

\begin{ejercicio}{2}{eje:t6-cuadratica}
Para los puntos $(-2,4),(-1,1),(0,0),(1,1),(2,4)$, en los que $y=x^{2}$: (a) calcule $r$ y la recta de mínimos cuadrados; (b) calcule los residuos y describa el patrón del gráfico de residuos; (c) ¿qué se concluye?
\end{ejercicio}

\begin{ejercicio}{3}{eje:t6-normales}
Demuestre el \cref{teo:t6-mc} por otra vía: (a) obtenga las ecuaciones normales igualando a cero las derivadas parciales de $\operatorname{SC}(b_0,b_1)$; (b) resuélvalas; (c) compruebe que la matriz hessiana es definida positiva, de modo que el punto crítico es un mínimo (y, por ser $\operatorname{SC}$ un polinomio de grado~2, global).
\end{ejercicio}

\begin{ejercicio}{3}{eje:t6-residuos}
Demuestre la \cref{prop:t6-residuos}: $\sum_ie_i=0$, $\sum_ix_ie_i=0$ y $\sum_i\widehat y_ie_i=0$. Deduzca que la media de los valores ajustados es $\ybar$.
\end{ejercicio}

\begin{ejercicio}{3}{eje:t6-origen}
(a) Demuestre que la recta que pasa por el origen y minimiza $\sum_i(y_i-bx_i)^{2}$ tiene pendiente $\widehat b=\sum_ix_iy_i\big/\sum_ix_i^{2}$ (suponga que algún $x_i\ne0$). (b) Demuestre que sus residuos cumplen $\sum_ix_ie_i=0$ y dé un ejemplo en el que $\sum_ie_i\neq0$.
\end{ejercicio}

\begin{ejercicio}{3}{eje:t6-transformaciones}
Sean $x_i'=\alpha x_i+\beta$ e $y_i'=\gamma y_i+\delta$ con $\alpha,\gamma\neq0$. Demuestre que la pendiente de mínimos cuadrados de $y'$ sobre $x'$ es $\widehat b_1'=(\gamma/\alpha)\,\widehat b_1$, que $R^{2}$ no cambia y que los residuos se multiplican por $\gamma$. ¿Cómo cambia la pendiente al pasar $x$ de metros a centímetros?
\end{ejercicio}

\begin{ejercicio}{3}{eje:t6-regresion-media}
Sean $z_{x,i}=(x_i-\xbar)/s_x$ y $z_{y,i}=(y_i-\ybar)/s_y$. (a) Demuestre que la recta de mínimos cuadrados, escrita en unidades tipificadas, es $\widehat z_y=r\,z_x$. (b) Deduzca que, si $\abs r<1$, la predicción de $y$ está a menos desviaciones típicas de su media que $x$ de la suya («regresión a la media»). (c) Un modelo obtiene en una primera evaluación una puntuación $2$ desviaciones típicas por encima de la media de los modelos; si la correlación entre dos evaluaciones es $0{,}5$, ¿qué puntuación tipificada se predice para la segunda?
\end{ejercicio}
